\documentclass[a4paper]{article}
\usepackage{amsfonts}
\usepackage{amsmath}
\usepackage{amssymb}
\usepackage{graphicx}     

\begin{document}
  
\noindent \large Auteur: Abdellah El Moussaoui \\
\noindent \large Studierichting: Informatica\\
\noindent \large Oefeningenreeks 5B nr. 3.2\\

\medskip

\normalsize

$\displaystyle\int^1_0\dfrac{\ln{x}}{\left(x+1\right)^2}dx$\\

$=\lim\limits_{b\to0+}\displaystyle\int^1_b\ln{x}\left(x+1\right)^{-2}dx$\\

We weten dat: $\displaystyle\int^c_dudv=\left[uv\right]^c_d-\displaystyle\int^c_dvdu$\\

Laat:\\

$u=\ln{x}$ , $dv=\left(x+1\right)^{-2}dx$\\

$du=\dfrac{1}{x}dx$ , $v=-\left(x+1\right)^{-1}$\\

Dus:\\

$\lim\limits_{b\to0+}\displaystyle\int^1_b\ln{x}\left(x+1\right)^{-2}dx = \lim\limits_{b\to0+} \left(-\left[\dfrac{\ln{x}}{x+1}\right]^1_b+\displaystyle\int^1_b\dfrac{1}{x\left(x+1\right)}dx\right)$\\

We weten dat:\\

$\dfrac{1}{x\left(x+1\right)} = \dfrac{A}{x}+\dfrac{B}{x+1}$\\

$\Rightarrow A=1$ , $B=-1$\\

$\lim\limits_{b\to0+}\displaystyle\int^1_b\ln{x}\left(x+1\right)^{-2}dx = \lim\limits_{b\to0+} \left(-\left[\dfrac{\ln{x}}{x+1}\right]^1_b+\displaystyle\int^1_b\left(\dfrac{1}{x}-\dfrac{1}{\left(x+1\right)}\right)dx\right)$\\

$= \lim\limits_{b\to0+} \left(-\left(\dfrac{\ln{1}}{2}-\dfrac{\ln{b}}{b+1}\right)+\left[\ln{x} - \ln{x+1} \right]^1_b\right)$\\

$= \lim\limits_{b\to0+} \dfrac{\ln{b}}{b+1}+\ln{1} - \ln2 -\ln{b} + \ln\left({b+1}\right)$\\

$= \lim\limits_{b\to0+} \dfrac{\ln{b}}{b+1} -\ln{b} + \ln\left({b+1}\right)- \ln2 $\\

$= \lim\limits_{b\to0+} \ln{b} \left(\dfrac{1}{b+1}-1\right) +\ln\left({b+1}\right)- \ln2 $\\


$= \lim\limits_{b\to0+} \ln{b} \left(\dfrac{-b}{b+1}\right) +\ln\left({b+1}\right)- \ln2 $\\

$ = 0 + 0 -\ln2$\\

$=-\ln2$\\


\begin{thebibliography}{999}
%\bibitem{a} Referentie
\end{thebibliography}
\end{document} 
